\documentclass[11pt]{amsart}
\ifdefined\pdfinfoomitdate\pdfinfoomitdate=1\relax\pdftrailerid{}\fi
\usepackage[T1]{fontenc}
\usepackage{lmodern,amsmath,amssymb,amsthm,mathtools,booktabs}
\usepackage[expansion=false]{microtype}
\usepackage{tikz}
\raggedbottom
\setlength{\emergencystretch}{2em}
\usepackage[margin=1in]{geometry}
\usepackage[numbers,sort&compress]{natbib}
\usepackage[hidelinks,hypertexnames=true]{hyperref}
\urlstyle{same}
\numberwithin{equation}{section}
\newtheorem{theorem}{Theorem}[section]
\newtheorem{proposition}[theorem]{Proposition}
\newtheorem{lemma}[theorem]{Lemma}
\newtheorem{corollary}[theorem]{Corollary}
\theoremstyle{definition}
\newtheorem{definition}[theorem]{Definition}
\theoremstyle{remark}
\newtheorem{remark}[theorem]{Remark}
\newcommand{\dd}{\,\mathrm{d}}
\newcommand{\eps}{\varepsilon}
\newcommand{\sgn}{\operatorname{sgn}}
\newcommand{\supp}{\operatorname{supp}}
\newcommand{\LieM}{\operatorname{Lie}_{M}}
\newcommand{\Tr}{\mathcal{T}}
\newcommand{\oscD}{\mathcal{D}}
\setcounter{tocdepth}{1}

\hypersetup{pdfauthor={Haibo Lu},pdftitle={Electronic supplement: Limit cycles on small supports},
  pdfsubject={Small-support constructions for Positive valleys and hyperbolic limit cycles in classical Lienard systems with a k-term primitive}}
\title[Electronic supplement: limit cycles on small supports]{Electronic supplement\\ Limit cycles on small supports}
\author{Haibo Lu}
\address{Shanghai Institute of Technology}
\date{November 27, 2026}
\thanks{ORCID: \href{https://orcid.org/0009-0000-2717-5968}{0009-0000-2717-5968}.}

\begin{document}
\maketitle
\noindent Electronic supplement to \emph{Positive valleys and hyperbolic limit
cycles in classical Li\'enard systems with a $k$-term primitive}. References prefixed
by \emph{Article} use the numbering of the main text.
\section{What a small support can realize}\label{sec:supp-intro}\hypertarget{public.audit}{}
We count limit cycles of a single system
\begin{equation}\label{eq:lienard}
  \dot{x}=y-F(x),\qquad \dot{y}=-x,\qquad F(0)=0,
\end{equation}
with real collected primitive
\begin{equation}\label{eq:collected}
  F(x)=\sum_{i=1}^k a_i x^{m_i},\qquad 1\le m_1<\cdots<m_k,\qquad a_i\ne0 .
\end{equation}
The exponents are integers and the degree is free. Let $H(k)$ and
$\LieM(k)$ be the respective suprema for hyperbolic and isolated cycles,
with values in $\mathbb Z_{\ge0}\cup\{+\infty\}$.
At six and seven terms a distant maximum closes ordinary and incoming
returns on the same axis. At eight and nine terms the rise and fall of
a single crest give two sets of sign brackets. Sections~3--6 prove the
bounds below. Section~7 identifies the finite inputs.

\begin{theorem}\label{thm:lower}
\leavevmode
\begin{enumerate}
  \item[\textup{(i)}] Chen, Jin, Li and Zhou attribute the one- and
  two-term counts for isolated cycles to Cen, Chen, Tian and Zhao
  \cite[p.~2]{CJLZ}. Section~\ref{subsec:baselines-lift}
  proves $H(1)=\LieM(1)=0$ and $H(2)=\LieM(2)=1$, including
  hyperbolic attainment in the two-term case.
  \item[\textup{(ii)}] For each of the following supports there is a
  nonempty open set in $(\mathbb R\setminus\{0\})^k$ on which the given
  number is a lower bound for the number of distinct hyperbolic limit cycles.
  In the last row every sufficiently large even $M$ works after the
  preceding profile has been fixed.
\[
\begin{array}{@{}rcl@{}}
\toprule
 k & \text{Cycles} & \supp(F)\\
\midrule
 3 & 3  & \{1,d-1,d\},\quad d\ge10\ \text{even}\\
 4 & 4  & \{3,5,9,10\}\\
 5 & 7  & \{5,9,10,29,16384\}\\
 6 & 10 & \{5,9,10,11,13,18\}\\
 7 & 11 & \{3,5,9,10,11,13,18\}\\
 8 & 14 & \{3,5,8,10,11,12,37,2^{38}\}\\
 9 & 19 & \{5,9,10,1033,2058,3081,3082,3083,M\}\\
\bottomrule
\end{array}
\]
  Six cycles also occur on
  $\{5,9,10,11,12\}$, and eleven on $\{5,9,10,11,13,15,M\}$ for every
  sufficiently large even $M$ in this second construction.
\end{enumerate}
\end{theorem}

\hypertarget{trinomial.extension}{}
\begin{corollary}[Odd-power extension of the three-cycle trinomial]\label{cor:odd-lift}
For every $k\ge3$, an open set of coefficients on
$\{1,9,10\}\cup\{2i+3:4\le i\le k\}$ realizes at least $k$
hyperbolic cycles.
\end{corollary}

Article, Proposition~A.1, proves this extension. Section~\ref{subsec:nine-powers}
applies it to a nineteen-cycle nest to obtain at least twenty cycles
on ten terms.

\subsection{Comparison on fixed supports}\label{subsec:full-support}
On $\{1,9,10\}$ the three simple return matches of Chen, Jin, Li and
Zhou give three hyperbolic cycles \cite[(3.1) and proof of
Theorem~1.2]{CJLZ}. Section~\ref{sec:trinomials} recovers this case
at $d=10$. Article, Proposition~A.1, successively adjoins the powers
$11,13,\ldots,21$, giving column two of Table~\ref{tab:low-comparison}.
At four terms Theorem~\ref{thm:lower}(ii) keeps the degree at ten.

For column three, the published count in degrees $3,4$ is exactly one
\cite{LMP,LL}. In degree $5$ the classical construction gives two
hyperbolic cycles \cite[p.~336]{LMP}. For $k\ge6$ the bound is $k-2$
in degree $k$ \cite[Theorem~1]{DMH}. By hyperbolic persistence these
examples can occupy the full support $\{1,\ldots,k\}$.

The unsigned Palomar preprint gives a rational sextic on full support
and a coefficient box of radius $10^{-40}$ in each coordinate. At
seven increasing positive heights its enclosures of both half-returns
give the signs $+,-,+,-,+,-,+$ for the difference $A^+-A^-$ of their
depths \cite[table \emph{Certified half-returns} and appendix]{Palomar}.
In our convention $b_\pm=-A^\pm$, so the signs reverse. Trapping between
the two bounding returns on each side supplies transverse returns
throughout the six brackets \cite[subsection \emph{Trapping}]{Palomar}.
Continuity allows their shared strict endpoints to be separated.

At each new degree the proof of \cite[Theorem~2]{Palomar} adjoins a
small nonzero term $cx^{n+1}$ while keeping the old coefficients and
strict samples. Each new odd degree supplies one more bracket by an
exterior test. For the specified septic $F_6+10^{-41}x^7$ this gives
seven. The full supports in degrees $6$ through $10$ thus carry
$6,7,7,8,8$ brackets. Palomar obtains isolated cycles throughout
the sextic box and for the specified septic. For column four, vary
the nonzero linear coefficient slightly. Its derivative in $b_+-b_-$
is positive on each return interval, so Article, Lemma~2.3 gives one
regular value with a hyperbolic cycle in every bracket. These cycles
persist on an open set.

Starting instead with the six nearby hyperbolic cycles and successively
adjoining $p=7,9,11,13$ by Article, Proposition~A.1, gives $8,9,10$
cycles on $8,9,10$ terms in degrees $9,11,13$. The
support $\{1,2,3,4,5,6,7,9,11,13\}$ of degree thirteen appears
in Article, Table~1. Lu's open exact-degree examples give column five
on full support in degrees $6$ through $9$, and at least fourteen
hyperbolic cycles on $\{1,\ldots,10\}$
\cite[supplement, Proposition~S1.2]{LuDegree}.

\begin{table}[htbp]
\centering\small
\setlength{\tabcolsep}{3.5pt}
\caption{Lower bounds for hyperbolic cycles on $k$ nonzero terms. Full
support means $\{1,\ldots,k\}$. Column~2 (`Seed or lift', from
\cite{CJLZ} and Article, Proposition~A.1) uses degrees
$10,11,13,15,17,19,21$. Column~3 gives published bounds
\cite{LMP,LL,DMH}. Column~4 uses the Palomar brackets at nearby
coefficients. Columns~5--6 use Lu's preprint \cite{LuDegree}, at full
support and on the indicated degree-twelve supports before closure.}
\label{tab:low-comparison}
\begin{tabular}{@{}rccccccr@{}}
\toprule
 & & \multicolumn{3}{c}{Full support, degree $k$}
 & Degree $12$ & \multicolumn{2}{c}{Theorem~\ref{thm:lower}(ii)}\\
\cmidrule(lr){3-5}\cmidrule(lr){6-6}\cmidrule(l){7-8}
$k$ & Seed or lift & \cite{LMP,LL,DMH} & \cite{Palomar} + Art.\ Lem.~2.3 & \cite{LuDegree}
    & Precl.\ \cite{LuDegree} & Cycles & Degree\\
\midrule
3 & 3 & 1 & -- & -- & -- & 3 & even $d\ge10$\\
4 & 4 & 1 & -- & -- & -- & 4 & 10\\
5 & 5 & 2 & -- & -- & -- & 7 & 16384\\
6 & 6 & 4 & 6 & 6 & -- & 10 & 18\\
7 & 7 & 5 & 7 & 7 & -- & 11 & 18\\
8 & 8 & 6 & 7 & 9 & 10 & 14 & $2^{38}$\\
9 & 9 & 7 & 8 & 10 & 11 & 19 & large even $M$\\
\bottomrule
\end{tabular}
\end{table}

Column~6 counts the ordinary returns before exterior closure in
Lu's degree-twelve construction \cite[\S\S5.2--5.3]{LuDegree}.
For one separated radius
($h=1$), Lu's Lemma~5.3 uses even powers $2,8,10,12$ and odd powers
$3,5,7$, with $b_7>0$. The linear term of his Proposition~3.1 gives the
allowed support $S_8=\{1,2,3,5,7,8,10,12\}$. In Lu's depth convention a
small launch gives $+$ and the eight tests near the radius alternate
from $-$ to $+$. By \cite[(3.8)]{LuDegree} a farther test has sign
$\operatorname{sgn}((7-12+2)b_7)=-$. Detuning the higher matched
reference to $+$ gives $1+7+1+1=10$ ordinary brackets.

For $S_9=S_8\cup\{9\}$ first add a sufficiently small negative
ninth-power coefficient $b_9$ to the profile. The preceding strict
tests persist and a new farther test has sign
$\operatorname{sgn}((9-12+2)b_9)=+$ \cite[Lemma~3.2]{LuDegree}.
Choose the reference above this test and detune it to $-$, giving
$1+7+2+1=11$ brackets. In each case one regular choice of the linear
coefficient gives a simple zero in every bracket
\cite[Proposition~2.4(i)]{LuDegree}. Hyperbolic persistence allows
any missing coefficients within $S_8$ or $S_9$ to be filled, giving
exactly eight or nine terms, respectively.

For the comparison per primitive term in Article, Section~1.3, stop
Lu's construction after $s\ge1$ separated radii in the proof of
\cite[Lemma~5.3]{LuDegree}. With $m=8+4s$ the available even powers
are $2,8,10,\ldots,m$ and the odd powers are $3,5,\ldots,5+2s$.
The ordinary, positive-height incoming and negative-height incoming
comparisons give at least $7s,3s,s$ changes. A small term of degree
$m-1$ preserves these signs. A linear coefficient realizes the
ordinary returns and regularizes their zeros while preserving the
incoming signs \cite[Lemma~3.2(ii) and Propositions~3.1, 2.4(i)]{LuDegree}.
At zero translation the two incoming sets of heights give at least
$4s$ changes. The exterior primitive adds at most the powers
$m+1,m+2,m+3$ and gives one additional cycle
\cite[Proposition~3.3, (3.9) and (3.11)]{LuDegree}. This proves the
count $7s+4s+1=11s+1$ on at most $(2s+2)+(s+2)+1+1+3=3s+9$
powers, in degree at most $m+3=4s+11$.

Keeping all the remote odd powers at one radius gives the sharper
finite counts used in the article. Set $h=1$ in the table in the
proof of \cite[Proposition~5.1]{LuDegree}. Its rows $e=0,1$ use
the same even powers $2,8,10,12$. Their highest inner odd powers
are $9$ and $11$, respectively. After adjoining the linear and
exterior powers the supports are
$\{1,\ldots,14\}\setminus\{4,6\}$ and
$\{1,\ldots,15\}\setminus\{4,6\}$. The counts $14h+9,14h+10$
therefore give $H(12)\ge23$ in degree $14$ and $H(13)\ge24$ in
degree $15$. With full supports in degrees $12,13,14$, the same
preprint gives $19,22,25$, respectively \cite[supplement, Theorem~S1.1
and Proposition~S1.2, and article, Theorem~1.1]{LuDegree}.

\subsection{Which tests a crest can transfer}
For a three-power seed we first match the two limiting lower intercepts
while retaining strict ordinary return brackets. The six-term construction
adds three powers at a remote fold. Three equations on finite half-returns
close incoming trajectories from prescribed heights in the same field
as the ordinary tests. The low-degree seven-term example then adds a
lower odd power. At five, eight and nine terms we use a different
transfer that requires two finite matches and no prescribed shift for
a list of incoming heights.

For the latter construction take a rising even primitive $E_b$,
$E_b'(r)=rf_b(r)>0$ for $r>0$. Subtract $f_b(1)r^M/M$ with $M$ even,
above every old even power and $M-2>\sup_{r>0}rf_b'(r)/f_b(r)$.
The resulting $E$ has one simple maximum at $1$. Choose an odd profile
$Q$ with $Q(1)=0$ and top power below $M$. With $f=E'/r$ put
$U(t)=\int_0^t(fQ'-f'Q)\dd r$. A test at $a<1$ and its partner
$z>1$ with $E(z)=E(a)$ read
\[
 U(a),\qquad U(a)-\frac{f(a)}{f(z)}[U(z)-U(1)] .
\]
To use $n$ changes on both branches, the two expressions must have the
same strict alternating signs at the $n+1$ chosen nodes. The mark
and the two far depths must also have the signs specified in
Proposition~\ref{prop:two-marks}. This is a comparison along the
rise and descent of the crest. Each test already compares the two
reflected half-planes. In the nine-term case seven brackets on each
branch, two simple marks, an exterior bracket and two bounded brackets
give nineteen. In the eight-term case the explicit wall exponent is
$2^{38}=274\,877\,906\,944$.

\section{Return conventions and the first two counts}\label{sec:toolkit}
In Article, Section~2 the left half-plane is reflected and time is reversed.
At a common upper launch $\beta>0$ with two transverse returns,
$b_\sigma(\beta)$ is the lower intercept for $P_\sigma(r)=F(\sigma r)$,
$r=|x|$. Zeros of $\oscD=b_+-b_-$ give cycles; at such a zero the
multiplier is $\exp(J_--J_+)$, where $J_\sigma=\int P_\sigma'(r)\dd t$
on the half-orbit. The comparison of ordered profiles assumes order on
the union of the two radial ranges and strict order on an interval that
both arcs traverse (Article, Lemma~2.1). The tail gap is
$\Delta=\ell_+-\ell_-$, $\ell_\sigma=\lim_{\beta\to\infty}b_\sigma(\beta)$.
The entering graph with limiting ordinate $h$ has axis ordinate $G_\sigma(h)$.

For radius and ordinate scales $\Lambda,M$, put $u=(r/\Lambda)^2/2$,
$w=y/M$, $\eps=(\Lambda/M)^2$ and $\phi_\sigma=P_\sigma(r)/M$.
The outgoing and incoming legs end at a common turn $T$ and have
positive and negative axis ordinates, respectively. Write $X,Y$ for
half their squared axis ordinates and $m_{\sigma,T},l_{\sigma,T}$ for
the suffix maximum and minimum of $\phi_\sigma$ on $[u,T]$.
Article (2.5)--(2.8) give the launch, coefficient and turn derivatives
and the envelope errors. Article, Lemma~2.2 gives the finite-turn
comparison and the incoming graph. Article, Lemma~2.3 supplies one
regular perturbation of an existing odd coefficient for any finite
list of strict brackets.
\subsection{One and two terms}\label{subsec:baselines-lift}

With two odd powers, an energy comparison orders any two putative cycles.
For attainment it suffices to perturb a single family of circles.

\begin{proof}[Proof of Theorem~\ref{thm:lower}\textup{(i)}]
An even $F$ has identical half-profiles. Every periodic orbit belongs to
an open annulus of periodic orbits by transversality of its two axis hits. This includes two even powers.
For a single odd monomial the divergence $-F'(x)$ has one strict sign off
$x=0$, so Green's theorem excludes a periodic orbit. This proves the
one-term assertion. With one odd and
one even power, $F(r)-F(-r)$ has a strict sign on $r>0$. A periodic
orbit would supply two transverse returns from its upper intercept.
Article, Lemma~2.1(ii) orders their lower hits strictly.

It remains to bound the number for two odd powers. Equal coefficient signs
again give divergence of one sign. With opposite signs we may reverse
time and replace $y$ by $-y$ so that
\[
 F(x)=-c x^{2p+1}+e x^{2q+1},\qquad c,e>0,\quad 0\le p<q .
\]
If $p=0$, the damping $F'$ is negative at zero and increases on
the positive half-axis, so Massera's uniqueness theorem applies
\cite[pp.~367--369]{Massera}. For $p>0$, $F'$ initially decreases.
The energy comparison below covers both cases. Let $a=(c/e)^{1/(2q-2p)}$. Then $F<0$ on $(0,a)$ and $F$ increases
strictly from zero on $[a,\infty)$. Central reflection preserves the
field's orbits. Each closed orbit is centrally symmetric, since its image
is either itself or a strictly nested closed orbit, and applying the
reflection twice rules out strict nesting.

Write $V=(x^2+y^2)/2$. Along an orbit $\dot V=-xF(x)$ and
$\dd V=F(x)\dd y$. On the right half of a closed orbit $V$ takes the
same value at the two axis hits. This arc must pass $x=a$, since $V$
strictly increases while $0<x<a$. Denote the upper and lower ordinates
at abscissa $x$ by $y_o(x)>F(x)>y_i(x)$. The gain of $V$ in $0<x<a$ is
\[
 \int_0^a x[-F(x)]\!
 \left(\frac1{y_o(x)-F(x)}+\frac1{F(x)-y_i(x)}\right)\dd x .
\]
For a strictly enclosing orbit the gain is smaller, since both vertical
gaps in the denominators increase. The loss beyond $a$ is
$\int_{y_i(a)}^{y_o(a)} F(x(y))\dd y$, where the right-hand arc is a
single graph over $y$. On the enclosing arc the interval of integration
is larger and $x(y)$ is strictly larger on the common interval.
As $F$ is strictly increasing there, the loss is larger. Two nested
cycles cannot both balance their gain and loss. This proves
$\LieM(2)\le1$, including $p>0$.

For attainment with a hyperbolic cycle take $F(x)=\delta(x^3-x)$.
On a compact annulus about radius $R_*=2/\sqrt3$, the return after one
clockwise revolution is smooth in $R$ and $\delta$. Along the unperturbed
circle $x=R\cos t$, integration gives
\[
 V_{\rm return}-V_{\rm initial}
  =\delta\!\int_0^{2\pi}(R^2\cos^2t-R^4\cos^4t)\dd t+O(\delta^2)
  =\pi\delta R^2(1-3R^2/4)+O(\delta^2)
\]
in $C^1$ on that annulus. Divide by $\delta$ and apply the implicit
function theorem at the simple positive zero $R_*$. For all sufficiently
small $\delta>0$ the resulting orbit is hyperbolic, proving $H(2)=1$.
\end{proof}

\section{Three hyperbolic cycles in every even degree \texorpdfstring{$d\ge 10$ ($k=3$)}{d >= 10 (k = 3)}}\label{sec:trinomials}

Chen, Jin, Li and Zhou construct three hyperbolic cycles for
$C(-Ax+Bx^9+x^{10})$ by varying $A$ along a curve of matched
half-returns \cite[\S3, Lemmas~3.1--3.2 and Theorem~1.2]{CJLZ}.
We use their matching curve and their comparison at two equal minima.
To treat every even $d\ge10$, we normalize the coefficient of $x^{d-1}$
inside the parentheses to one. A test height at which both turn
equations factor then gives an action gap with a uniform sign for
$d\ge20$ and rational enclosures for the five smaller even degrees.

Work first with $F(x)=C(x^d+x^{d-1}-ax)$, where $d\ge 10$ is even and $a,C>0$, so the scaled half-profiles are
\begin{equation}\label{eq:trinomial-profiles}
  f_\sigma(x,a) = x^d + \sigma(x^{d-1} - ax) \qquad (x\ge 0,\ \sigma=\pm 1).
\end{equation}
Set $\tau_*=(d-2)/(d-1)$, $\rho_*=(d-1)/d$, $a_c=\tau_*^{d-2}/(d-1)=\max_{[0,1]}x^{d-2}(1-x)$, and
$K_*=\rho_*^{d-1}/d$. For $0<a<a_c$, $f_+'(x,a)=dx^{d-1}+(d-1)x^{d-2}-a$ has a unique positive root
$s\in(0,1)$, where $f_+$ attains its global minimum, while $f_-'(x,a)=x^{d-2}(dx-(d-1))+a$ first
decreases and then increases on $(0,\infty)$, taking a negative value at $\tau_*$ and the positive
value $a$ at $\rho_*$, so $f_-$ attains its negative global minimum at its outer critical point
$r\in(\tau_*,\rho_*)$. Both minima are nondegenerate, and $\partial_a[f_+(s,a)-f_-(r,a)]=-s-r<0$.
Since $(f_+(s,a),f_-(r,a))$ equals $(0,-K_*)$ at $a=0$ and $(f_+(s,a_c)<0,\,0)$ at $a=a_c$, there is a
unique $a_d\in(0,a_c)$ at which $f_+(s,a_d)=f_-(r,a_d)<0$.

In the degree-ten proof, the level $60$ after scaling is chosen above
the inner maxima \cite[Lemma~3.2, (3.14)--(3.19)]{CJLZ}.
Fix $a=a_d$ and choose the test height $H=a_d>0$. The exact factorizations
\begin{equation}\label{eq:trinomial-fact}
  f_+(x,a) - a = (x+1)(x^{d-1}-a), \qquad f_-(x,a) - a = (x-1)(x^{d-1}+a)
\end{equation}
show that $f_+$ and $f_-$ stay strictly below $H=a_d$ on $[0,t)$ and $[0,1)$, respectively, and cross
$H$ transversely at $t=a_d^{1/(d-1)}<1$ and $1$. Define the equal-depth contraction difference
\begin{equation}\label{eq:Id-def}
  I_R = \int_s^t \frac{f_+'(x,a_d)^2}{x}\dd x, \qquad I_L = \int_r^1 \frac{f_-'(x,a_d)^2}{x}\dd x, \qquad I_d = I_R - I_L.
\end{equation}

\begin{lemma}[Uniform negative contraction gap for $d\ge 10$]\label{lem:trinomial-Id}
For every even integer $d\ge 10$, one has $I_d<0$; for every even $d\ge 20$, $d I_d < -71/6480$.
\end{lemma}

\begin{proof}
At $a=a_d$, comparing the equal minima via $f_+(s,a)\ge -a$ and $f_-(r,a)\le
f_-(\rho_*,a)=-K_*+a\rho_*$ gives $a\ge K_*/(1+\rho_*)$. Using $f_+(s,a)\le
f_+(2/3,a)=(5/3)(2/3)^{d-1}-(2/3)a$ and $f_-(r,a)=r^{d-1}(r-1)+ar\ge -K_*+ar$ gives $ar-K_*\le
(5/3)(2/3)^{d-1}-(2/3)a$, so $(r+2/3)a\le K_*+(5/3)(2/3)^{d-1}$ with $r>\tau_*=(d-2)/(d-1)$.

Assume $d\ge 20$. Since $(1-1/d)^{d-1}=(1+1/(d-1))^{-(d-1)}>1/e>1/3$ and $(19/20)^{19}<19/50$, we have
$1/3<dK_*<19/50$. Both $d(5/3)(2/3)^{d-1}$ and $d(3/4)^{d-1}$ decrease for $d\ge 20$, and at $d=20$
are strictly less than $1/60$ and $1/6$, respectively. Hence
\[
  \frac{1}{6d} = \frac{1/3}{2d} < a_d \le \frac{dK_* + d(5/3)(2/3)^{d-1}}{d(\tau_* + 2/3)} <
  \frac{19/50 + 1/60}{18/19 + 2/3}\,\frac{1}{d} < \frac{1}{4d},
\]
and $(3/4)^{d-1}<1/(6d)<a_d=t^{d-1}$ implies $t>3/4$. On $[s,t]$ we have $0\le f_+'(x,a_d)\le
(dt+d-1)x^{d-2}$. Extending the integral to $[0,t]$ and using $t^{d-1}=a_d<1/(4d)$ and $t>3/4$ gives
\[
  d I_R \le \frac{d}{2(d-2)}(da_d)^2\left(1+\frac{d-1}{dt}\right)^2 <
  \frac{20}{36}\cdot\frac{1}{16}\cdot\left(\frac{7}{3}\right)^2 = \frac{245}{1296}.
\]
On $[\rho_*,1]\subset[r,1]$, dropping $+a_d>0$ from $f_-'(x,a_d)=x^{d-2}(dx-(d-1))+a_d>0$ and setting
$w=d(1-x)\in[0,1]$ (since $2d-5<2d$ and $(1-w/d)^d\ge 1-w$ by Bernoulli's inequality on $[0,1]$)
yields
\[
  d I_L \ge \int_0^1 \left(1-\frac{w}{d}\right)^{2d-5}(1-w)^2\dd w \ge \int_0^1 (1-w)^4\dd w = \frac{1}{5}.
\]
Thus $dI_d < 245/1296 - 1/5 = -71/6480 < 0$ for all even $d\ge 20$. For $d=10,12,14,16,18$, exact
rational interval arithmetic (Table~\ref{tab:trinomial-cert} in Section~\ref{sec:reproduce}) certifies
$1000dI_d\in(-24,-22)$, $(-51,-50)$, $(-69,-67)$, $(-81,-79)$, and $(-89,-88)$, respectively.
\end{proof}

\begin{lemma}[Global matching curve for $F(x)=C(x^d+x^{d-1}-ax)$]\label{lem:trinomial-curve}
Fix even $d\ge 10$ and $C>0$.
\begin{enumerate}
  \item[\textup{(i)}] For every $\beta>0$ there exists a unique $a=\mu_C(\beta)>0$ at which both
  half-arcs from $(0,\beta)$ return transversely to the same negative intercept
  $b_+(\beta,a)=b_-(\beta,a)<0$. The function $\mu_C:(0,\infty)\to(0,a_c+2/C)$ is real analytic,
  satisfies $\lim_{\beta\downarrow 0}\mu_C(\beta)=0$, and along $a=\mu_C(\beta)$ one has
  $\sgn\mu_C'(\beta)=\sgn(J_+(\beta)-J_-(\beta))$ and $\Pi'(\beta)=\exp(J_-(\beta)-J_+(\beta))$.
  \item[\textup{(ii)}] For fixed $d,C$, $\mu_C'(\beta)>0$ for all sufficiently large $\beta$.
  \item[\textup{(iii)}] As $C\to+\infty$, $\mu_C(Ca_d)\to a_d$ and $C^{-2}\bigl(J_+(Ca_d)-J_-(Ca_d)\bigr)\to I_d$.
\end{enumerate}
\end{lemma}

\begin{proof}
We adapt Chen, Jin, Li and Zhou's degree-ten matching curve
\cite[Lemmas~3.1--3.2]{CJLZ}. The degree-dependent checks are as follows.

(i) At $a=0$ coercive tails and rotation at the origin give both
transverse returns, and $f_+>f_-$ on $x>0$ gives $\oscD>0$.
For $a>0$ the plus origin repels and its tail forces a turn.
On the common return domain Article~(2.6) gives
$\partial_a b_+<0<\partial_a b_-$. Lowering $a$ lowers the minus
profile: its graph in decreasing $y$ stays outside any returning
reference arc, is bounded by coercivity and crosses $y=0$ at positive
radius. After that it hits the axis since $\dot x=y-Cf_-(x,a)$
is negative near $x=0$ at a negative ordinate. Thus the minus return
domain at fixed $\beta$ is an interval beginning at zero.
For $a\ge a_c+2/C$ we have $Cf_-(x,a)\ge2x$; the wedge
$0\le x\le y$ is invariant. At the upper endpoint of the return
interval the bounded minus arc approaches the hyperbolic sink.
A trapping neighborhood gives $b_-\to0$ there while $b_+<0$.
Strict monotonicity gives the unique analytic match. For
$0<a<2/C$ the two linear focus half-turns give
$\oscD(\beta,a)/\beta\to-2\sinh(\pi Ca/\sqrt{4-C^2a^2})<0$
as $\beta\downarrow0$, hence $\mu_C(\beta)\to0$.
Article~(2.5) gives the derivative sign and multiplier at a match.

(ii) For $\beta\to\infty$ put $R=(\beta/C)^{1/d}$, let
$x_\sigma=\sqrt{2T_\sigma}$ be the turning radius, and write
$G(u)=Cf_\sigma(\sqrt{2u},a)$ with $a=\mu_C(\beta)$ bounded.
Article~(2.8) gives $\beta=G(T_\sigma)+O(x_\sigma)$ and hence
$x_\sigma=R-\sigma/d+O(R^{-1})$. The common lower hit $\alpha$
is bounded away from zero by $b_+(1,0)<0$ for $\beta\ge1$.
Choose $K$ with $G(u)\ge u-K$; the barrier $u-y-K-2<0$
persists from the turn to the axis since its $y$-derivative is at
least one at equality. Thus $\alpha\ge-K-2$.

On the inward leg the gap $z=G-w_i$ solves $z'=G'-1/z$ with
$z(T_\sigma)=0$. Beyond a fixed $U$ both $G',G''>0$; a crossing
of $z=1/G'$ would have positive relative derivative, so $0<z<1/G'$
throughout that strip. In $\int G'/z$ use $1/z=G'-z'$ and integrate
by parts: the remainder involves $\int G''z\le\log(G'(T_\sigma)/G'(U))$
and bounded endpoint terms. On $[0,U]$ the inward gap stays bounded
below: at zero use $-\alpha$, and on a regular compact strip a small
gap forces $(z^2)'\le-1$ until it closes before the distant turn.
The slope is integrable at zero. The outgoing leg contributes
$O(x_\sigma^{d-1})$: on $[U,T_\sigma]$ use
$|G'|=O(x_\sigma^{d-2})$ and the ordinate change at most
$x_\sigma$; on $[0,U]$ its gap grows with $\beta$.
Consequently
\[
 J_\sigma=C^2\left(\frac{d^2 x_\sigma^{2d-2}}{2d-2}
 +\sigma\frac{2d(d-1)x_\sigma^{2d-3}}{2d-3}\right)+O_{d,C}(R^{2d-4}).
\]
The turn shifts contribute $-2d$ to the leading difference. Hence
$J_+-J_-\sim [2d/(2d-3)]C^2R^{2d-3}>0$.

(iii) At the fixed height $H=a_d$, \eqref{eq:trinomial-fact}
places the two outer roots on regular rising branches above all earlier
maxima, uniformly near $a=a_d$. With $\Lambda=1,M=C$, Article~(2.8)
gives convergence of $b_\sigma(CH,a)/C$ to the respective minima.
Their difference crosses zero with derivative $-s-r$, so
$\mu_C(CH)\to a_d$. The actual turns converge uniformly to the outer
roots. On a short initial interval the two legs have gaps bounded away
from zero and $\int|\phi_\sigma'|\dd u<\infty$. On the remaining
interval apply Article~(2.7)--(2.8) and integrate by parts. The
outgoing envelope is constant; the incoming envelope follows the
profile from its minimum to the turn. Thus $C^{-2}J_{o,\sigma}\to0$
and $C^{-2}J_{i,\sigma}$ tends to the corresponding integral of
$(f_\sigma')^2/x$, uniformly near $a_d$. Evaluating at the matching
parameter gives \eqref{eq:Id-def}.
\end{proof}

\begin{proof}[Proof of the $k=3$ row of Theorem~\ref{thm:lower}\textup{(ii)}]
As in the proof of Theorem~1.2 of \cite{CJLZ}, a regular horizontal
level of the matching curve yields three simple crossings.
Fix even $d\ge 10$. By Lemmas~\ref{lem:trinomial-Id} and~\ref{lem:trinomial-curve}(i),(iii), fix $C>0$
so large that $\mu_C'(Ca_d)<0$, and choose $0<\beta_1<\beta_2$ near $Ca_d$ with
$\mu_C(\beta_1)>\mu_C(\beta_2)$. By Lemma~\ref{lem:trinomial-curve}(ii), $\mu_C$ attains its minimum
on $[\beta_2,\infty)$ at some finite $\beta_m\ge\beta_2$ and admits $\beta_3>\beta_m$ with
$\mu_C(\beta_3)>\mu_C(\beta_m)$. Because $\mu_C$ is real analytic and nonconstant, its critical values
are at most countable; choose a regular value
$a\in(\mu_C(\beta_m),\min\{\mu_C(\beta_1),\mu_C(\beta_3)\})$. Since $\lim_{\beta\downarrow
0}\mu_C(\beta)=0<a$, the graph of $\mu_C$ crosses level $a$ upward in $(0,\beta_1)$, downward in
$(\beta_1,\beta_m)$, and upward in $(\beta_m,\beta_3)$, yielding three distinct hyperbolic limit
cycles (attracting, repelling, attracting). Rescaling $(x,y)=C^{-1/(d-1)}(X,Y)$ transforms
$C(x^d+x^{d-1}-ax)$ into $X^d+C^{1/(d-1)}X^{d-1}-CaX$, and hyperbolic persistence yields a nonempty
open neighborhood of $(A,b)=(Ca,C^{1/(d-1)})\in(0,\infty)^2$ on $\{1,d-1,d\}$.
\end{proof}


\section{A balanced trinomial and its first extensions}\label{sec:balanced}\hypertarget{seed.balance}{}
For the primitive $x^{10}+bx^9+\lambda x^5$, Article, Lemma~3.1
gives the balancing interval and the compact brackets. Its proof
encloses the two turning radii and the action difference at each of
the heights $20,100,200$.
\subsection{Four powers in degree ten and five powers in degree twelve}\label{subsec:four-five}

The same degree-ten trinomial supports two extensions. The first inserts
a cycle near the origin. The second preserves three finite cycles and
closes three more outside them. For both we first match the half-returns at a
finite launch height, where the derivative of their difference can be computed.

\begin{proof}[Proof of the four-term and degree-twelve assertions of Theorem~\ref{thm:lower}]
Use $F_{C,a}(x)=C(x^{10}+b_0x^9-ax^5)$ from
Article, Lemma~3.1, and write $S_\pm(h)$ for the two
squared-slope integrals there. The heights $20,100,200$ exceed the
inner crest. At $200C$, $b_\pm(200C)/C$ tend to the respective
minimum values of $q_{\pm,a}$. These values cross strictly at $a_0$
as $a$ varies. Bracketing $a_0$ on either side, then using continuity
and strict odd-coefficient comparison, gives a unique nearby
$\tilde a_C\to a_0$ with $\oscD(200C)=0$. This is a finite-height
condition, distinct from equality of the two tail limits.

Set $K_\sigma(h)=-\dd[b_\sigma(Ch)]/\dd h>0$ at this coefficient.
On $20\le h\le200$, the running-minimum estimate for
Article~(2.5) gives $-C^{-2}\log K_\sigma(h)\to S_\sigma(h)$
uniformly. The two actions increase strictly, and the prefactors have
logarithms $o(C^2)$ since $b_\sigma(Ch)/C$ tends to the common
negative minimum. For $h<200$ we have
\[
 \oscD(Ch)=\int_h^{200}(K_+(s)-K_-(s))\dd s.
\]
The integral of each positive $K_\sigma$ has rate $S_\sigma(h)$,
attained at the left endpoint. Thus the smaller rate determines the
sign. At the mark itself one compares $-K_+(200)+K_-(200)$.
The rational enclosures for $S_+-S_-$ at $20,100,200$ are
$(-89,-88)$, $(1293,1295)$, and $(6884,6886)$, respectively
(Section~\ref{sec:reproduce}). Hence $\oscD(20C)>0$, $\oscD(100C)<0$,
and $\oscD'(200C)>0$. The negative lowest odd term gives
$\oscD(\beta_0)<0$ at a small $\beta_0>0$.

Fix $C$ and add $+\gamma x^3$ with $\gamma>0$ small. The three
strict signs and the simple zero persist. On a still smaller orbit
segment the cubic dominates, giving $\oscD(\beta_1)>0$ for some
$0<\beta_1<\beta_0$. The three resulting sign brackets and the marked
orbit give four cycles on $\{3,5,9,10\}$ after regularization by
Article, Lemma~2.3.

For the other extension, keep the cubic coefficient zero and first
regularize the two old brackets and the mark using the coefficient of
$x^5$. Scale both coordinates by the same positive constant to make
the primitive monic. Within this three-cycle open set, the tail gap
$\Delta=\ell_+-\ell_-$ can be made nonzero by varying that same
coefficient. Reflect $x$ and reverse time if necessary to make $\Delta>0$.
Fix a height $B_0$ above the three cycles with $\oscD_P(\beta)>\Delta/2$
for $\beta\ge B_0$, and put
\[
 F_{R,\kappa}(x)=P(x)-\kappa R^{-1}x^{11}+R^{-2}x^{12}.
\]
We now match at $2R^{10}$. For $\kappa=2-\delta$, $0<\delta<1$,
the bound $1-\kappa r+r^2\ge1-\kappa^2/4$ gives a uniform positive
tail barrier on both halves. Article, Lemma~2.2 gives
$\oscD(2R^{10})\to\Delta$. For $\kappa=2+\delta$, the right arc
passes $R$. Otherwise its outward envelope would give
$2\le\max_{0\le r\le1}r^{10}(1-r)^2+o(1)<1/100+o(1)$.
On its inward leg the ordinate at $R$ is below
$F_{R,\kappa}(R)=(-\delta+o(1))R^{10}$. Its lower-axis hit tends
to $-\infty$, whereas the other hit tends to $\ell_-(P)$. The sign
therefore reverses. All powers are at least five, so these half-arcs
return for every $\kappa$ under consideration. Scalar comparison gives
a unique matching $\kappa_R\to2$.

With $r=x/R$, $w=y/R^{10}$, the limiting profiles are
$p_\pm(r)=r^{10}(1\mp r)^2$. For $1/100\le h\le2$, let $t_\pm(h)$
be the inverse on the last increasing branch and set
\[
 S_\pm^{\rm out}(h)=\int_{\rho_\pm}^{t_\pm(h)}
             p_\pm'(r)^2\,\frac{\dd r}{r},\qquad \rho_+=1,\quad\rho_-=0 .
\]
The lower limits follow from the running minima. The two lower hits
stay in a compact subinterval of $(-\infty,0)$, between their returns
at $B_0$ and the common return at $2R^{10}$, which tends to $\ell_-(P)$.
Since the slopes $p_\pm'/r$ of the rescaled profiles have uniformly
bounded variation down to zero, the envelope estimate and
Article~(2.5) give these actions as the uniform rates of
$-\dd[b_\pm(R^{10}h)]/\dd h$ at speed $R^{18}$. At $h=1/100,1,2$,
the three enclosures for $S_+^{\rm out}-S_-^{\rm out}$ are
$(3/10000,4/10000)$, $(-3/2,-1)$, and $(-5,-4)$.
Integrating back from the mark as above gives
\[
 \oscD(B_0)>0,\qquad \oscD(R^{10}/100)<0,\qquad \oscD(R^{10})>0,
 \qquad \oscD(2R^{10})=0,\quad \oscD'(2R^{10})<0 .
\]
Compact convergence preserves the three original hyperbolic cycles.
Two new brackets and the simple mark lie above $B_0$. Regularization
in the $x^5$ coefficient gives six simultaneous hyperbolic cycles
on $\{5,9,10,11,12\}$.
\end{proof}


\section{Ten cycles in degree eighteen}\label{sec:six}
\subsection{Matching at a distant fold}\label{subsec:supp-fold}
Apply Article, Proposition~4.4 and Corollary~4.5, with $d_{\rm top}=10$
and $m_0=5$. The three new powers are $11,13,18$. Write
$\beta_{\rm refl}$ for the upper launch at the end of the descending
fold bracket and $\beta_{\rm ext}$ for the plus upper hit at the
exterior test depth $s_0<S_*$. Between them lie the launches
corresponding to the prescribed shifted incoming tests, in decreasing
order of physical turning height. Both endpoint displacements are
positive. The two simple marks and both fold brackets lie outside
this open interval.
\subsection{Ten cycles on six powers and eleven cycles on seven powers in degree eighteen}\label{subsec:six-power-proof}

The balanced trinomial already has two compact brackets and two
incoming tests. The fold of Article, Proposition~4.4 closes
the incoming trajectories in degree eighteen. The second construction
uses the empty interval between the origin and the innermost of these
cycles to add one more power.

\begin{proof}[Proof of the $k=6$ row of Theorem~\ref{thm:lower}\textup{(ii)}]
Take the monic balanced degree-ten seed $P_{\mathrm{seed},\lambda}(x)=x^{10}+bx^9+\lambda x^5$
($d_{\mathrm{top}}=10$, $m_0=5$) from Article, Lemma~3.1(ii)--(iii). It provides
$N_{\mathrm{ord}}=2$ bounded ordinary sign brackets and two negative heights
$0>h_{\mathrm{seed},1}>h_{\mathrm{seed},2}$ at which $G_+(h)-G_-(h)$ at $\lambda_0$ has signs $-,+$.
Fix $s>0$ so small that $G_+(h_{\mathrm{seed},1})-G_-(h_{\mathrm{seed},1}+s)<0$ and
$G_+(h_{\mathrm{seed},2})-G_-(h_{\mathrm{seed},2}+s)>0$, and choose
$h_{\mathrm{deep}}<h_{\mathrm{seed},2}$ so negative that $|G_\pm(h)-h|<s/4$ for all $h\le
h_{\mathrm{deep}}+s$ (by Article, Lemma~2.2(iv)), ensuring
$G_+(h_{\mathrm{deep}})-G_-(h_{\mathrm{deep}}+s)<-s/2<0$. Apply Article, Proposition~4.4 to
$(h_1,h_2,h_3)=(h_{\mathrm{seed},1},h_{\mathrm{seed},2},h_{\mathrm{deep}})$ with signs $(-,+,-)$.
From $\beta_{\rm refl}$ to $\beta_{\rm ext}$ the signs $(+,-,+,-,+)$
give four disjoint brackets. With the two bounded seed brackets, two
fold brackets and two marked cycles, Article, Proposition~4.4
gives ten simultaneous hyperbolic limit cycles on $\{5,9,10\}\cup\{11,13,18\}=\{5,9,10,11,13,18\}$.
\end{proof}

\begin{corollary}[Explicit degree-eighteen seven-power support]\label{cor:seven-inner}
A nonempty open subset of $(\mathbb{R}\setminus\{0\})^7$ on the degree-$18$ support
$\{3,5,9,10,11,13,18\}$ realizes at least $11$ simultaneous hyperbolic limit cycles.
\end{corollary}

\begin{proof}
Fix a six-term polynomial $F_0$ on $\{5,9,10,11,13,18\}$ with ten hyperbolic limit cycles, and isolate
their positive $y$-axis intercepts in pairwise disjoint compact intervals. Below the innermost
interval, the linearized rotation gives two transverse returns for
small launches. Choose $\beta_0>0$ so both arcs lie in a neighborhood of
$(0,0)$ where the odd part of $F_0$ has the strict sign of its lowest odd coefficient $a_5\ne 0$; by
Article, Lemma~2.1(ii), $\oscD_{F_0}(\beta_0)$ has the strict sign $\sgn(a_5)$. Adjoin
$\gamma x^3$ of opposite sign $\sgn(\gamma)=-\sgn(a_5)$ with $|\gamma|>0$ small enough that all ten
hyperbolic cycles and the strict sign at $\beta_0$ persist. At a still smaller height
$\beta_1\in(0,\beta_0)$, the cubic term $\gamma r^3$ dominates the odd part on both half-arcs, giving
$\sgn\oscD(\beta_1)=\sgn(\gamma)=-\sgn\oscD(\beta_0)$ and bracketing an eleventh periodic orbit in
$(\beta_1,\beta_0)$. Applying Article, Lemma~2.3 to the $x^3$ coefficient makes all eleven
orbits simultaneously hyperbolic while keeping all seven powers in $\{3,5,9,10,11,13,18\}$ nonzero.
\end{proof}

\section{Two matched orbits at a single crest}\label{sec:two-matches}
A high even power can bend an increasing primitive over a crest. The
ascending and descending turns then see the same height at different
radii. Two exact matches fix the constants of integration in the return
comparison. A list on the ascent gives a second set of brackets only
when the descending variation, including motion of the paired turn,
has the required strict signs at the same limiting levels.
Proposition~\ref{prop:two-marks} specifies the sign at each node and
the additional conditions at the marked turn and at the two outer depths.

\subsection{Two finite-return matches across a single crest}\label{subsec:two-marks}
Let $E_b(r)$ be an even polynomial with $E_b(0)=0$ and $E_b'(r)=r f_b(r)$, where $f_b(r)>0$ for all
$r>0$ and $f_b$ has lowest term $d r^{d-2}$ ($d\ge 8$ even). Choose an even integer $M$ strictly
exceeding every exponent of $E_b$ and satisfying $M-2>\sup_{r>0} r f_b'(r)/f_b(r)$, and set
\begin{equation}\label{eq:crest-wall}
  E(r) = E_b(r) - \frac{f_b(1)}{M}r^M, \qquad f(r) = \frac{E'(r)}{r} = f_b(r) - f_b(1)r^{M-2}, \qquad H_c = E(1) > 0.
\end{equation}
Because $f_b(r)/r^{M-2}$ decreases strictly on $(0,\infty)$, $f(r)$ is positive on $(0,1)$, vanishes
simply at $r=1$, and is negative on $(1,\infty)$, so $E$ attains a unique positive critical point at
$r=1$ (a strict global maximum $H_c$). For $y<H_c$, let $z(y)>1$ denote the unique descending root of
$E(z)=y$, and put $g=z(0)>1$. Let $Q(r)$ be an odd polynomial with lowest term $r^q$ ($3<q<d$ odd),
all remaining powers strictly above $d$, highest odd power $P<M$ carrying a nonzero coefficient
$[r^P]Q\ne 0$, and $Q(1)=0$. Define
\begin{align}
  U(a) &= \int_0^a \bigl(f(r)Q'(r)-f'(r)Q(r)\bigr)\dd r, \qquad \Tr(t) = U(t) - U(1), \nonumber\\
  B_L(a) &= U(a), \qquad B_R(a) = U(a) - \frac{f(a)}{f\bigl(z(E(a))\bigr)}\Tr\bigl(z(E(a))\bigr). \label{eq:two-actions}
\end{align}
The correction in $B_R$ comes from matching the upper hits of
descending turns, as in the preprint formula
\cite[supplement, p.~16, (S5.12)--(S5.13)]{LuDegree}; Article,
Lemma~4.2\textup{(ii)}, gives a self-contained proof of the common
leading rate of the two derivative terms through the crest and the
first-order formulas $B_L,B_R$ uniformly on compact sets of regular
turns.

Combine $(E,Q)$ with a three-power inner seed $P_{\mathrm{seed},\lambda}(x)=x^d+bx^q+\lambda x^l$
($b>0$, $3\le l<q$ odd) on a compact interval $I_\lambda\subset(-\infty,0)$ across whose endpoints the
tail gap $\Delta(P_{\mathrm{seed},\lambda})=\ell_+-\ell_-$ has opposite strict signs while $N$
pairwise disjoint bounded ordinary brackets persist throughout $I_\lambda$:
\begin{itemize}
  \item When $(l,q,d)=(5,9,10)$, Article, Lemma~3.1 supplies $I_\lambda\subset(-\infty,0)$ and $N=2$ bounded brackets.
  \item When $(l,q,d)=(3,5,8)$ with $b=1$, we take $N=0$ and verify that an interval
  $I_\lambda\subset(-\infty,0)$ straddles tail balance: at $\lambda=0$,
  $P_{\mathrm{seed},0}(r)=r^8+r^5>r^8-r^5=P_{\mathrm{seed},0}(-r)$ on $r>0$, so
  $\Delta(P_{\mathrm{seed},0})>0$ by Article, Lemma~2.2(iii); for $\lambda<0$, evaluating an
  inward graph at $r=1$ gives
  $\ell_+(P_{\mathrm{seed},\lambda})<P_{\mathrm{seed},\lambda}(1)=2+\lambda$, whereas monotonicity in
  $\lambda$ (Article, Lemma~2.2(iii)) gives
  $\ell_-(P_{\mathrm{seed},\lambda})>\ell_-(P_{\mathrm{seed},0})>-\infty$, so
  $\Delta(P_{\mathrm{seed},\lambda_1})<0$ for $\lambda_1<\ell_-(P_{\mathrm{seed},0})-3$ and continuity
  yields $\Delta(P_{\mathrm{seed},\lambda_2})>0$ for $\lambda_2\in(\lambda_1,0)$ close to $0$.
\end{itemize}
For large radial scale $R>1$ and $|\theta|\le\delta:=\frac{1}{2}|[r^P]Q|>0$, define the physical primitive
\begin{equation}\label{eq:two-field}
  F_{R,\lambda,\theta}(x) = R^d E(x/R) + b R^q\bigl[Q(x/R) + \theta (x/R)^P\bigr] + \lambda x^l.
\end{equation}
Because $|\theta|<|[r^P]Q|$ and $\lambda\in I_\lambda$ avoids $0$, every listed power in
$\supp(F_{R,\lambda,\theta})$ carries a nonzero coefficient.

\begin{proposition}[Finite single-crest transfer]\label{prop:two-marks}
Fix $0<a_0<\dots<a_n<a_*<1$ and $\sigma\in\{+1,-1\}$. Suppose that for $0\le i\le n$, both $B_L(a_i)$
and $B_R(a_i)$ have strict sign $\sigma(-1)^i$, while $B_L(a_*)$, $\Tr(z(-1/2))$, and $\Tr(z(-1))$ all
have strict sign $\sigma(-1)^n$. Then for all sufficiently large $R>1$, there exist
$(\lambda_R,\theta_R)\in I_\lambda\times[-\delta,\delta]$ and an arbitrarily small perturbation of
$\lambda_R$ such that system \eqref{eq:lienard} for $F_{R,\lambda,\theta_R}$ has at least
\[
  N + 2n + 2 + \mathbf{1}_{\{n\text{ odd}\}} + \mathbf{1}_{\{N=0,\ \sigma=-1\}}
\]
simultaneous hyperbolic limit cycles on $\supp(F_{R,\lambda,\theta_R})$.
\end{proposition}

\begin{proof}
Pass to outer scaled coordinates $r=|x|/R$, $u=r^2/2$, $w=y/R^d$, $\tau=R^{q-d}$, and $\eps=R^{2-2d}$,
so the scaled half-profiles are $p_\pm(r)=E(r)\pm b\tau(Q(r)+\theta r^P)\pm\lambda R^{l-d}r^l$.
Because $q\ge 5>1$ and $l<q$, one has $\sqrt{\eps}=R^{1-d}=o(\tau)$, $\eps|\log\eps|=o(\tau)$, and
$R^{l-d}=o(\tau)$.

On the rectangle $I_\lambda\times[-\delta,\delta]$ we solve two simultaneous finite-return equations:
\begin{enumerate}
  \item At the upper intercept $\beta(R a_*)$ of the plus turn $r=Ra_*$, require
  $b_+(\beta(Ra_*))=b_-(\beta(Ra_*))$. Choose $t_0\in(a_*,1)$: on $[0,t_0]$, $E(r)/r^d$ has a positive
  minimum, so $F_{R,\lambda,\theta}(\pm x)\ge c x^d$ holds uniformly on $K_0\le x\le R t_0$, while
  $F_{R,\lambda,\theta}\to P_{\mathrm{seed},\lambda}$ in $C^1$ on compacts along converging parameters. The finite turns meet both axes
  by rotation near the origin. For $0<t_1<a_*<t_0$ their rescaled upper hits
  at $Rt_1,Ra_*,Rt_0$ tend to $E(t_1)<E(a_*)<E(t_0)$ for each sign by
  Article~(2.8). Thus $\beta(Ra_*)\to\infty$ and its minus
  turn also lies below $Rt_0$. Both inward legs start at the profile
  value at their own turn. The finite-turn comparison of Article, Lemma~2.2(ii)
  gives $b_+(\beta(Ra_*))-b_-(\beta(Ra_*))\to\Delta(P_{\mathrm{seed},\lambda})$ uniformly in
  $\theta\in[-\delta,\delta]$, yielding opposite strict signs on $\partial
  I_\lambda\times[-\delta,\delta]$.
  \item At common scaled lower depth $s=1$ (physical lower intercept $-R^d$, i.e., $Y_\pm(\widehat
  T_\pm(1))=1/2$ with $\sqrt{2\widehat T_\pm(1)}\to z(-1)>1$), require $X_+(\widehat
  T_+(1))=X_-(\widehat T_-(1))$. Because $z(-1)>1$ lies past the perturbed crests $\rho_\pm=1+O(\tau)$
  and $Q(1)=0$, Article~(2.8) encloses $\sqrt{2X_\pm(\widehat T_\pm(1))}$ within
  $O(\sqrt{\eps})$ of the peak heights $p_\pm(\rho_\pm)=H_c\pm b\tau\theta+O(\tau^2+R^{l-d})$, so
  $\sqrt{2X_+(\widehat T_+(1))}-\sqrt{2X_-(\widehat T_-(1))}=2b\tau\theta+o(\tau)$ has opposite strict
  signs on $I_\lambda\times\{\pm\delta\}$.
\end{enumerate}
By the two-dimensional Poincar\'e--Miranda theorem, for all large $R$ there exists a simultaneous
solution $(\lambda_R,\theta_R)\to(\lambda_0,0)$, at which the $N$ bounded seed brackets persist by
compact $C^1$ convergence.

To apply Article, Lemma~4.2, fix $T_{\max}>z(-1)^2/2$ and
$\mathcal I=[0,T_{\max}]$ after choosing the wall exponent $M$. The later
turns $\widehat T_\pm(1)$ lie inside $\mathcal I$ for large $R$ and the second
match gives equal upper hits there. At the matched parameters put
$\delta_{\rm c}=bR^{q-d}$. The remainders beyond
$E\pm\delta_{\rm c}Q$ have $C^1(\mathcal I)$ norm
$O(\delta_{\rm c}|\theta_R|+R^{l-d})=o(\delta_{\rm c})$ in $u$,
since $l\ge3$ and $l<q$. The slope of each monomial has bounded
variation on $\mathcal I$, uniformly in $R$, and away from $u=0$ the profiles
converge in $C^2$ to $E(\sqrt{2u})$, preserving the simple crest.
Moreover $\sqrt\eps/\delta_{\rm c}=R^{1-q}/b\to0$.
On the early ascent, division by $r^{d-1}$ in the radial derivative
bounds the two lower powers by multiples of $(Rr)^{q-d}$ and
$(Rr)^{l-d}$. Thus every additional critical point lies at
$r\le K_0/R$ for fixed sufficiently large $K_0$.
There the physical slopes divided by the physical radius are bounded,
so $\phi_\pm'=O(R^{2-d})$ on an interval of $u$-length $O(R^{-2})$.
Its action is $O(R^{2-2d})=o(\delta_{\rm c})$.
Together with the transverse hits this proves the hypotheses up to
the depth-one reference. The lemma now gives
$\sgn\oscD(\beta(Ra_i))=-\sigma(-1)^i$ on the ascent and
$\sgn\oscD(\beta(Rz(E(a_i))))=\sigma(-1)^i$ on the descent.
Since $X_+$ is strictly increasing, the descending tests occur in
reverse index order. There are $n$ strict brackets on each side and
a simple marked orbit at $a_*$. The last descending sign is $\sigma$.

For the two exterior depths take $v_0=1/2$, $V=1$ and
$\alpha=b\tau=\delta_{\rm c}$. The expansion $p_\pm=E\pm\alpha Q+o(\alpha)$
holds at the matched parameters with the uniform bounds just established.
Use the exterior comparison in the proof of Article, Proposition~4.4.
Here $z(-v)$ remains on the regular descent for $1/2\le v\le1$ and
both legs return transversely. At common lower depth $v$ the actual
turn levels are $-v+O(\sqrt\eps)$ by Article~(2.8). Therefore the
common cost $C(v)=\int_1^{z(-v)} f^2r\dd r$ increases strictly, and
the signed splitting there is $\pm\alpha\Tr(z(-v))+o(\alpha)$ in
$-\eps J_\pm$, uniformly. For $K(v)=(\beta_+(v)^2-\beta_-(v)^2)/(2R^{2d})$,
the upper match gives $K(1)=0$. The nonzero $\Tr(z(-1/2))$ fixes the
sign of $-K'$ near $1/2$; the cost gap in the cited integral estimate
makes that neighborhood dominate $K(1/2)=-\int_{1/2}^1K'(v)\dd v$.
The nonzero $\Tr(z(-1))$ gives $K'(1)\ne0$. Since $b_-'<0$, passing
from the upper-hit difference at fixed lower hit to $\oscD$ at the
plus launch preserves the sign and preserves simplicity at the mark.
At depth $1/2$ this gives $\sgn\oscD(\beta_+(1/2))
=\sgn\Tr(z(-1/2))=\sigma(-1)^n$. This turn has radius asymptotic to
$Rz(-1/2)$ and lies beyond $Rz(E(a_0))$. The second marked orbit
lies at depth $1$. When $n$ is odd, $\sigma$ and $\sigma(-1)^n=-\sigma$ are opposite, creating one
additional bracket between $Rz(E(a_0))$ and depth $1/2$. When $N=0$ and $\sigma=-1$, the first
ascending test has $\oscD(\beta(Ra_0))>0$ (at height $\asymp R^d E(a_0)\to+\infty$), while for fixed
$\lambda_R<0$ the lowest odd term $\lambda_R x^l$ forces $\oscD(\beta_0)<0$ at a small $\beta_0>0$,
creating one inner bracket in $(\beta_0,\beta(Ra_0))$. Slightly shrinking shared bracket endpoints and
applying Article, Lemma~2.3 to the $x^l$ coefficient makes all counted periodic orbits
simultaneously hyperbolic.
\end{proof}

\subsection{Five powers ($7$ cycles) and seven powers ($11$ cycles)}\label{subsec:five-seven}
Take the balanced degree-ten seed $(l,q,d)=(5,9,10)$ ($N=2$) with pure monomial slope $f_b(r)=10r^8$,
so for $Q(r)=\sum b_p r^p$ one has $U_{f_b,Q}(r)=10r^8\sum c(p)b_p r^p$ with $c(p)=(p-8)/(p+8)$.

\begin{proof}[Proof of the $k=5$ and $k=7$ rows of Theorem~\ref{thm:lower}\textup{(ii)}]
For $k=5$, set $Q(r)=r^9-r^{29}$ and fix the explicit wall exponent $M=16384$. Bisection with dyadic test points, outward interval evaluation on the
decreasing branch and polynomial antiderivatives give the enclosures below.
At $(a_0,a_1,a_*)=(\frac{5}{6},\frac{19}{20},\frac{39}{40})$, the scaled pairs
$1000\bigl(B_L(a)/f(a),\,B_R(a)/f(a)\bigr)$ lie in
\[
  (8,9)\times(2,3), \qquad (-92,-91)\times(-97,-96), \qquad (-226,-225)\times(-231,-230),
\]
while $\Tr(z(-1/2))\in(-176,-175)$ and $\Tr(z(-1))\in(-247,-246)$. Proposition~\ref{prop:two-marks}
with $N=2$, $n=1$, and $\sigma=+1$ yields $2+2+2+1=7$ simultaneous hyperbolic limit cycles on
$\{5,9,10,29,16384\}$.

For this five-term support let $z_i=z(E(a_i))$ and
write $\beta(t)$ here for the upper hit of the plus turn at physical
radius $Rt$. Let $\beta_v$ be its upper hit at scaled lower depth $v$.
At the simultaneous solution, the following are the actual
lower-intercept signs in increasing order of upper launch.
\[
\begin{array}{@{}lcccccc@{}}
\toprule
 &\text{seed}& \beta(a_0),\beta(a_1)&\beta(a_*)&
 \beta(z_1),\beta(z_0)&\beta_{1/2}&\beta_1\\
\midrule
\sgn(b_+-b_-)&-,+,-&-,+&0&-,+&-&0\\
\bottomrule
\end{array}
\]
The two zeros displayed are simple. The sign following $\beta(z_0)$
supplies the fifth strict bracket. The two marked orbits complete the count.

The same test applies to $Q=r^9-r^{11}$ and $M=20{,}000{,}000$.
At $a=1/2,3/4,4/5$ both actions have signs $+,-,-$; the two depth
actions lie in $(-39,-38)$ and $(-53,-52)$, respectively
(Section~\ref{sec:reproduce}). This also gives seven cycles on
$\{5,9,10,11,20{,}000{,}000\}$.

For $k=7$, put $p_j=9+2j$ ($0\le j\le 3$), $(\xi_1,\xi_2,\xi_3)=(\frac{1}{192},\frac{1}{96},1)$, and
\[
  V(z) = c(9)\prod_{j=1}^3\left(1 - \frac{z}{\xi_j}\right), \qquad b_{p_j} = \frac{[z^j]V}{c(p_j)},
  \qquad Q_{\mathrm{raw}}(r) = \sum_{j=0}^3 b_{p_j}r^{p_j}.
\]
Then $U_{f_b,Q_{\mathrm{raw}}}(r)=10r^{17}V(r^2)$ has three simple positive roots
$\sqrt{\xi_1}<\sqrt{\xi_2}<1$. Using $1/c(p)=1+16/(p-8)$ and $V(1)=0$, exact integration gives
$Q_{\mathrm{raw}}(1)=16\int_0^1 V(t^2)\dd t>0$, while $b_{15}<0$. Because $U_{f_b,Q_{\mathrm{raw}}}'$
has at least three positive roots by Rolle's theorem and at most three by Descartes' rule,
$(Q_{\mathrm{raw}}/f_b)'$ has strict negative sign on $(1,\infty)$, giving a unique simple root $C>1$
of $Q_{\mathrm{raw}}$ with $Q_{\mathrm{raw}}'(C)<0$. Rescaling $Q(r):=C^{-9}Q_{\mathrm{raw}}(Cr)$
places all three simple roots of $U_{f_b,Q}$ inside $(0,1)$, satisfies $Q(1)=0$ and
$Q'(1)=C^{-8}Q_{\mathrm{raw}}'(C)<0$, and allows us to fix $0<a_0<a_1<a_2<a_3<a_*<1$ at which
$U_{f_b,Q}$ alternates through signs $+,-,+,-,-$.

Now let $M>15$ be an even integer tending to $+\infty$. At each fixed $a_i$ and $a_*$ in $(0,1)$,
$B_L(a)\to U_{f_b,Q}(a)$. For any fixed regular height $y<E_b(1)$ (specifically
$y=E_b(a_i)+O(M^{-1})$, $y=-1/2$, and $y=-1$), solving $f_b(1)z^M=M(E_b(z)-y)$ gives $M\log z=\log
M+O(1)$, i.e., $z-1=(\log M)/M+O(M^{-1})$, while $f(z)\sim -M(E_b(1)-y)<0$ and $\int_1^z|f(r)|\dd r\le
H_c-y=O(1)$. Integration by parts therefore yields
\begin{equation}\label{eq:wall-correction}
  \frac{\Tr(z)}{f(z)} = -Q(z) + \frac{2\int_1^z f(r)Q'(r)\dd r}{f(z)} = -Q'(1)\frac{\log M}{M} + O(M^{-1}).
\end{equation}
As $M\to+\infty$, $\Tr(z(E(a_i)))/f(z(E(a_i)))\to 0$, so $B_R(a_i)$ retains the strict signs $+,-,+,-$
of $U_{f_b,Q}(a_i)$; simultaneously, since $f(z)<0$ and $Q'(1)<0$, multiplying
\eqref{eq:wall-correction} by $f(z)$ gives $\Tr(z(-s))\sim (E_b(1)+s)Q'(1)\log M<0$ at $s=1/2$ and
$s=1$ (sign $-1=\sigma(-1)^3$). Fixing any sufficiently large even $M$,
Proposition~\ref{prop:two-marks} with $N=2$, $n=3$, and $\sigma=+1$ yields $2+6+2+1=11$ simultaneous
hyperbolic limit cycles on $\{5,9,10,11,13,15,M\}$.
\end{proof}

\subsection{Eight powers at explicit degree $2^{38}$ ($14$ cycles)}\label{subsec:eight-powers}

For $k=8$ we work from the three-power seed $x^8+x^5+\lambda x^3$ ($d=8$, $q=5$, $l=3$, $b=1$, $N=0$)
and insert two intermediate even powers $\{10,12\}$ so that $f_b(r)=E_b'(r)/r>0$ develops a deep
positive dip without vanishing, amplifying the ascending action reversals of
$Q(r)=r^5-512r^{11}+511r^{37}$ to five before closing the crest with the explicit high even exponent
$M=2^{38}=274{,}877{,}906{,}944$.

\begin{proof}[Proof of the $k=8$ row of Theorem~\ref{thm:lower}\textup{(ii)}]
Set $d=8$, $q=5$, $l=3$, $z=8r^2$, and $C(z)=1-\frac{19}{10}z+z^2>0$. In \eqref{eq:crest-wall} define
\[
  f_b(r) = 8r^6 C(z), \qquad Q(r) = r^5 - 512r^{11} + 511r^{37}, \qquad M = 2^{38}.
\]
Here $Q(1)=0$, $Q'(1)=13280$, $A_1:=f_b(1)=1992/5$, and integrating $r f_b(r)$ gives
\[
  E_b(r) = r^8 - \frac{304}{25}r^{10} + \frac{128}{3}r^{12}, \qquad E_b(1) = \frac{2363}{75}.
\]
Using $C(z)\ge(1-\frac{19}{20})(1+z^2)=\frac{1}{20}(1+z^2)$, the logarithmic slope satisfies
$rf_b'(r)/f_b(r)=6+2zC'(z)/C(z)\le 6+4\cdot 20=86<M-2$, so the single-crest monotonicity condition
holds on $(0,\infty)$. At the seven ordered nodes
\[
  z \in \left\{\frac{1}{5},\, \frac{1}{2},\, 1,\, 2,\, \frac{11}{2},\, 7,\, \frac{15}{2}\right\}
\]
(the last serving as the marked ascending turn $a_*=\sqrt{15/16}<1$), writing $G(z)=1-z^3+511(z/8)^{16}$
expresses the unclosed ratio $U_{f_b,Q}(r)/f_b(r)$ in exact rational form up to $(z/8)^{5/2}$:
\[
  \frac{U_{f_b,Q}(r)}{f_b(r)} = \left(\frac{z}{8}\right)^{5/2}\frac{1}{C(z)}\sum_{i,j} C_i G_j\,\frac{-1+2(j-i)}{11+2(i+j)}\,z^{i+j}.
\]
Exact rational arithmetic gives the bounds in Section~\ref{sec:reproduce} and shows that these seven
values have signs $-,+,-,+,-,+,+$ and magnitudes exceeding $1/400{,}000$. Because $r^2\le 15/16$ at
all seven nodes, each integrated cross-term with $-A_1 r^{M-2}$ has exponent factor of modulus at most
$1$, $\sum|[r^p]Q|=1024$, and $f_b(r)\ge \frac{2}{5}r^6$, so subtracting the $M=2^{38}$ wall changes
$U/f_b$ by at most $2^{20}(15/16)^{(M-8)/2}<10^{-9}$. Since $0<f(r)\le f_b(r)$ at the test nodes,
$B_L/f=U_{f,Q}/f$ retains the signs $-,+,-,+,-,+,+$ with magnitudes exceeding $1/410{,}000$.

At $M=2^{38}$ we check the descent at the positive test levels and at
the two negative levels. All ascending test levels lie in $(0,21)$ because $E(r)$ increases on $(0,1)$ and
$E_b(\sqrt{15/16})<21$. For every $y\in[-1,21]$, the descending inverse $z(y)>1$ solving $A_1
z^M=M(E_b(z)-y)$ satisfies
\[
  1 + \frac{16}{M} < z(y) < 1 + \frac{40}{M}, \qquad -f\bigl(z(y)\bigr) > \frac{1}{2}M(H_c-y).
\]
Indeed, $(1+16/M)^M<3^{16}$ whereas $(1+40/M)^M\ge\exp(40-800/M)>2^{39}$; throughout $[1,1+40/M]$ one
has $r^{37}<4/3$, whence $E_b(r)<114$ and $f_b(r)<1300$ by triangle bounds on their coefficients,
while $E_b(r)\ge E_b(1)>31.5$ brackets $z(y)$ inside $(1+16/M,1+40/M)$. Furthermore, $1/z^2>3/4$,
$H_c=E(1)>31$, and $f(z)=f_b(z)-M(E_b(z)-y)/z^2$, proving $-f(z(y))>\frac{1}{2}M(H_c-y)$.

On $[1,1+40/M]$, $|Q''(r)|<2\cdot 736{,}992$ implies $13{,}000<Q'(r)<14{,}000$. Substituting these
bounds into the exact integration-by-parts identity
\[
  \frac{\Tr(z)}{f(z)} = -Q(z) + \frac{2\int_1^z f(r)Q'(r)\dd r}{f(z)}
\]
(using $f(1)=Q(1)=0$ and $\int_1^z|f(r)|\dd r\le\int_1^z|E'(r)|\dd r=H_c-y$) gives $Q(z)>16\cdot
13{,}000/M$ while the integral ratio has modulus at most $4\cdot 14{,}000/M$. Consequently
$\Tr(z)/f(z)<0$ (whence $\Tr(z)>0$ because $f(z)<0$, giving $\Tr(z(-1/2))>0$ and $\Tr(z(-1))>0$) and
\[
  \left|\frac{\Tr(z)}{f(z)}\right| \le \frac{(40+4)\cdot 14{,}000}{M} < \frac{1}{440{,}000} < \frac{1}{410{,}000}.
\]
Hence at each ascending test $a_i$ ($0\le i\le 5$), the descending expression
$B_R(a_i)/f(a_i)=B_L(a_i)/f(a_i)-\Tr(z(E(a_i)))/f(z(E(a_i)))$ preserves every strict sign
$-,+,-,+,-,+$. Applying Proposition~\ref{prop:two-marks} with $N=0$, $n=5$, and $\sigma=-1$ yields
$0+5+5+2+1+1=14$ simultaneous hyperbolic limit cycles on $\{3,5,8,10,11,12,37,2^{38}\}$.
\end{proof}


\subsection{Nine powers via rational cubic separation ($19$ cycles)}\label{subsec:nine-powers}
At $k=9$, starting from the balanced seed $\{5,9,10\}$ ($N=2$) we insert two intermediate even powers
$\{2058,3082\}$ and three odd powers $\{1033,3081,3083\}$ before closing the crest with a large even
wall $M>3083$, achieving $n=7$ ascending action sign changes and $2+14+2+1=19$ simultaneous hyperbolic
limit cycles.

\begin{proof}[Proof of the $k=9$ row of Theorem~\ref{thm:lower}\textup{(ii)}]
Start with $f_0(r)=10r^8$ and $Q_{\mathrm{base}}(r)=r^9$. Set $D_9=1024$, $h=2^{-30}$, $T_0=15/8$, and define the positive cubics
\[
  C(z) = (1-z)^2(1+2z) + h z^2, \qquad G(z) = \left(1-\frac{z}{T_0}\right)^2\left(1+\frac{z}{2T_0}\right) + h z.
\]
For $z=r^{D_9}$, put $f_b(r)=10r^8 C(z)$ (strictly positive on $r>0$, integrating to $E_b$ with even
primitive exponents $\{10,2058,3082\}$) and $Q_0(r)=r^9 G(z)$ (odd exponents $\{9,1033,3081\}$).
Termwise integration gives $U_{f_b,Q_0}(r)=10 r^{17} J_{17}(z)$, where
\[
  J_{17}(z) = \sum_{i,j} C_i G_j\,\frac{1 + D_9(j-i)}{17 + D_9(i+j)}\,z^{i+j}.
\]
At the seven ordered nodes
\[
  z \in \left\{\frac{1}{16},\, \frac{1}{6},\, 1,\, \frac{15}{8},\, 4,\, 32,\, 4096\right\},
\]
exact rational evaluation gives the signs $+,-,+,-,+,-,+$ for $J_{17}(z)$ with $|J_{17}(z)|>1/1000$ at every node.

Here the top exponents of $f_b$ and $Q_0$ are $e=3080$ and $p=3081$, both with positive leading
coefficients. To append a seventh sign reversal and create the crest zero of $Q$, set $P=3083$ and
define $Q_B(r)=Q_0(r)-Q_0(B)(r/B)^P$ for large $B>1$. As $B\to+\infty$, $Q_B\to Q_0$ on compact
intervals so the seven alternating signs $+,-,+,-,+,-,+$ persist at the fixed nodes above, while at
$r=B$ one has $Q_B(B)=0$, $Q_B'(B)<0$ (since $P>p$ and $Q_0$ has positive coefficients at $p=3081$ in
dominance), and dividing $U_{f_b,Q_B}(B)$ by the product of the positive leading coefficients of
$f_b,Q_0$ and $B^{e+p}$ gives the strict negative limit
\[
  \frac{p-e}{p+e} - \frac{P-e}{P+e} = \frac{1}{6161} - \frac{3}{6163} < 0.
\]
Fix $B$ sufficiently large, choose a seventh test $r_7<B$ and a marked turn $r_*<B$ near $B$ where
$U_{f_b,Q_B}<0$, and rescale $r\mapsto Br$ (dividing $E_b$ by $B^{10}$ and $Q_B$ by $B^9$ to preserve
the lowest seed normalization $r^{10}$ and $r^9$) so that the eight ascending points
$0<a_0<\dots<a_7<a_*<1$ carry signs $+,-,+,-,+,-,+,-,-$ and the rescaled odd polynomial satisfies
$Q(1)=0$ and $Q'(1)<0$.

Finally, choose any sufficiently large even integer $M>3083$ satisfying
$M-2>\sup_{r>0}rf_b'(r)/f_b(r)$ and such that the asymptotic bound \eqref{eq:wall-correction}
preserves all eight ascending signs on $B_L(a_i)$ and $B_R(a_i)$ ($i=0,\dots,7$), preserves
$B_L(a_*)<0$, and gives $\Tr(z(-1/2))<0$ and $\Tr(z(-1))<0$ (matching $\sgn Q'(1)=-1=\sigma(-1)^7$).
Proposition~\ref{prop:two-marks} with $N=2$, $n=7$, and $\sigma=+1$ then yields $2+14+2+1=19$
simultaneous hyperbolic limit cycles on $\{5,9,10,1033,2058,3081,3082,3083,M\}$.
\end{proof}

\smallskip\noindent\emph{Twenty cycles on ten terms.}
For every sufficiently large even $M>3083$ after the profiles above
are fixed, the nineteen-cycle nest is hyperbolic. Since $M+1$ is odd
and exceeds the degree, Article, Proposition~A.1, adjoins $x^{M+1}$
with an arbitrarily small nonzero coefficient of suitable sign.
The nineteen cycles persist and the new exterior cycle is hyperbolic,
giving at least twenty on this ten-term support.

\section{Finite inputs for the return comparisons}\label{sec:reproduce}\label{sec:finite-inputs}
\hypertarget{finite.inputs}{}

For the seed and fold tests below, the programs evaluate rational
polynomial antiderivatives at endpoints isolated by bisection.
The improper integrals in the ten-term examples are bounded using
outward-rounded real balls in Arb \cite{Johansson2017}.

\subsection{Trinomials and the two extensions in low degree}\label{subsec:finite-seed}
Table~\ref{tab:trinomial-cert} supplies the finite cases in
Lemma~\ref{lem:trinomial-Id}. Its analytic estimate gives
$dI_d<-71/6480$ for the remaining even $d\ge20$. At $d=6,8$ the same equal-minimum integrals give $I_d>0$.
The negative action difference used for $d\ge10$ occurs in neither
of these two rows.

\begin{table}[htbp]
  \centering
  \small
  \caption{Exact rational enclosures $L<1000dI_d<U$ on $a_d\in[c-10^{-9},c+10^{-9}]$ for the trinomial
  equal-depth contraction difference $I_d$ in \eqref{eq:Id-def}.}
  \label{tab:trinomial-cert}
  \begin{tabular}{@{}rrrr@{}}
    \toprule
    $d$ & $c$ & $L$ & $U$ \\
    \midrule
    $6$  & $0.062500000$ & $130$ & $135$ \\
    $8$  & $0.040486856$ & $25$  & $27$  \\
    $10$ & $0.029550858$ & $-24$ & $-22$ \\
    $12$ & $0.023104700$ & $-51$ & $-50$ \\
    $14$ & $0.018889538$ & $-69$ & $-67$ \\
    $16$ & $0.015933561$ & $-81$ & $-79$ \\
    $18$ & $0.013753521$ & $-89$ & $-88$ \\
    \bottomrule
  \end{tabular}
\end{table}

For the degree-ten seed, the equal-minimum radius, peak, two
curvatures and action gaps at $20,100,200$ have rational enclosures
in the proof of Article, Lemma~3.1. Section~\ref{subsec:four-five}
combines those three actions with the tests at $1/100,1,2$ on the
outer degree-twelve profile.

\subsection{The fold and the single crest}\label{subsec:finite-crest}
The six-power construction uses the degree-eighteen fold data, and
Corollary~\ref{cor:seven-inner} adds a lower odd power. The five-power
examples require rising, reflected and depth tests for two walls;
the alternative seven-power construction starts with a cubic.
For eight powers the rising-node margin is $1/410000$ and the wall
exponent $2^{38}$ bounds the reflected correction by $1/440000$
(Section~\ref{subsec:eight-powers}). For nine powers, the seven rational
inputs to $J_{17}$ have alternating signs and absolute values greater
than $1/1000$. The new odd power gives one further sign by
$1/6161-3/6163<0$. Proposition~\ref{prop:two-marks} uses these tests
after the final even exponent is fixed.

\subsection{Reproducing the calculations}\label{subsec:replay}
In the standalone computation archive, \texttt{README.txt} lists each
test with its command and the file containing its reference output.
The standard suite covers the rational seed enclosures, the
degree-eighteen fold and the tests in Proposition~\ref{prop:two-marks}.
For Article, Corollary~4.6, it encloses the seed and positive-height
variations at $768$ bits and the five depth integrals with analytic
tails. A central difference at $\pm2^{-240}$ also checks the ordinary
variation at the moving minimum and its upstream inverse. Article,
Lemma~3.4 proves the signs for all four prefix choices. The suite
checks the fixed degree-$26$ profile at $512$ bits and refines the
integrals in degrees $26$ and $67088$. A separate $512$-bit check uses
reciprocal substitution for degree $22$.
The archived enclosures concern the local inputs to Article,
Corollary~4.6 and Corollary~5.2.

\begingroup
\raggedright
\renewcommand{\bibfont}{\small}
\providecommand{\doi}[1]{DOI \nolinkurl{#1}}
\bibliographystyle{abbrvnat}
\bibliography{tex/references}
\endgroup

\end{document}
